%TCIDATA{Version=5.50.0.2890}
%TCIDATA{LaTeXparent=0,0,NoticeDKE.tex}

\subsubsection{Density}

From calculations in Sec.\ref{moments}, the flux surface averaged density
$\left\langle n_{e}\right\rangle _{V,l+1/2}$ at $\psi_{l+1/2}$ may be
expressed as a sum of three terms
\[
\left\langle n_{e}\right\rangle _{V,l+1/2}=\left\langle n_{e}\right\rangle
_{V,l+1/2}^{0}+\left\langle n_{e}\right\rangle _{V,l+1/2}^{1}+\left\langle
\widetilde{n}_{e}\right\rangle _{V,l+1/2}^{1}
\]
the first one corresponding to the zero order Fokker-Planck equation
\begin{equation}
\left\langle n_{e}\right\rangle _{V,l+1/2}^{0}=2\pi\frac{\widetilde{q}%
_{l+1/2}}{\widehat{q}_{l+1/2}}\sum_{i=0}^{n_{p}-1}\sum_{j=0}^{n_{\xi_{0}}%
-1}p_{i+1/2}^{2}\lambda^{l+1/2,j+1/2}f_{0,l+1/2,i+1/2,j+1/2}^{(0)}\Delta
p_{i+1/2}\Delta\xi_{0,j+1/2}%
\end{equation}
while the two others result from the solution of the electron drift kinetic
equation
\begin{equation}
\left\langle n_{e}\right\rangle _{V,l+1/2}^{1}=2\pi\frac{\widetilde{q}%
_{l+1/2}}{\widehat{q}_{l+1/2}}\sum_{i=0}^{n_{p}-1}\sum_{j=0}^{n_{\xi_{0}}%
-1}p_{i+1/2}^{2}\lambda^{l+1/2,j+1/2}g_{l+1/2,i+1/2,j+1/2}^{(0)}\Delta
p_{i+1/2}\Delta\xi_{0,j+1/2}%
\end{equation}
and%

\begin{equation}
\left\langle \widetilde{n}_{e}\right\rangle _{V,l+1/2}^{1}=2\pi\frac
{\widetilde{q}_{l+1/2}}{\widehat{q}_{l+1/2}}\sum_{i=0}^{n_{p}-1}\sum
_{j=0}^{n_{\xi_{0}}-1}p_{i+1/2}^{2}\overline{\lambda}_{1,-1,0}^{l+1/2,j+1/2}%
\widetilde{f}_{l+1/2,i+1/2,j+1/2}^{(0)}\Delta p_{i+1/2}\Delta\xi_{0,j+1/2}%
\end{equation}

\subsubsection{Current Density}

In a similar way, the flux surface averaged parallel current $\left\langle
J_{\Vert}\right\rangle _{\phi,l+1/2}$ at $\psi_{l+1/2}$ may be expressed as a
sum of three terms
\begin{equation}
\left\langle J_{\Vert}\right\rangle _{\phi,l+1/2}=\left\langle J_{\Vert
}\right\rangle _{\phi,l+1/2}^{0}+\left\langle J_{\Vert}\right\rangle
_{\phi,l+1/2}^{1}+\left\langle \widetilde{J}_{\Vert}\right\rangle
_{\phi,l+1/2}^{1}%
\end{equation}
where the zero order term is
\begin{align}
&\left\langle J_{\Vert}\right\rangle _{\phi,l+1/2}^{0}=\frac{2\pi q_{e}}{m_{e}%
}\frac{q_{l+1/2}}{\overline{q}_{l+1/2}}\sum_{i=0}^{n_{p}-1}\sum_{j=0}%
^{n_{\xi_{0}}-1}\frac{p_{i+1/2}^{3}}{\gamma_{i+1/2}}H\left(  \left|
\xi_{0,j+1/2}\right|  -\xi_{0T,l+1/2}\right)\nonumber\\&\quad\times  \xi_{0,j+1/2}%
f_{0,l+1/2,i+1/2,j+1/2}^{(0)}\Delta p_{i+1/2}\Delta\xi_{0,j+1/2}%
\end{align}

The first order term arising from function $g$ is
\begin{align}
&\left\langle J_{\Vert}\right\rangle _{\phi,l+1/2}^{1}=\frac{2\pi q_{e}}{m_{e}%
}\frac{q_{l+1/2}}{\overline{q}_{l+1/2}}\sum_{i=0}^{n_{p}-1}\sum_{j=0}%
^{n_{\xi_{0}}-1}\frac{p_{i+1/2}^{3}}{\gamma_{i+1/2}}H\left(  \left|
\xi_{0,j+1/2}\right|  -\xi_{0T,l+1/2}\right) \nonumber\\&\quad\times \xi_{0,j+1/2}%
g_{l+1/2,i+1/2,j+1/2}^{(0)}\Delta p_{i+1/2}\Delta\xi_{0,j+1/2}%
\end{align}
while the other one which results from $\widetilde{f}$ is more
complex and
\begin{align}
&\left\langle \widetilde{J}_{\Vert}\right\rangle
_{\phi,l+1/2}^{1}=\frac{2\pi
q_{e}}{m_{e}}\frac{\widetilde{q}_{l+1/2}}{\overline{q}_{l+1/2}}\frac{R_{p}%
}{R_{0,l+1/2}}\frac{B_{T0,l+1/2}}{B_{0,l+1/2}}\sum_{i=0}^{n_{p}-1}\sum
_{j=0}^{n_{\xi_{0}}-1}\frac{p_{i+1/2}^{3}}{\gamma_{i+1/2}}\lambda
_{2,-2,2}^{l+1/2,j+1/2}\nonumber\\&\quad\times\xi_{0,j+1/2}\widetilde{f}_{l+1/2,i+1/2,j+1/2}%
^{(0)}\Delta p_{i+1/2}\Delta\xi_{0,j+1/2}.
\end{align}

\subsubsection{Power Density Associated with a Flux}

The density of power absorbed by the plasma through a particular mechanism is
the sum
\[
\left\langle P_{abs}^{\mathcal{O}}\right\rangle _{V,l+1/2}^{0}=\left\langle
P_{abs}^{\mathcal{O}}\right\rangle _{V,l+1/2}^{0}+\left\langle P_{abs}%
^{\mathcal{O}}\right\rangle _{V,l+1/2}^{1}+\left\langle \widetilde{P}%
_{abs}^{\mathcal{O}}\right\rangle _{V,l+1/2}^{1}%
\]
where the respective contributions are given by the equations (\ref{pabs0})
and (\ref{pabs1t}-\ref{pabs1g}) and discretized as%
\begin{align}
\left\langle P_{abs}^{\mathcal{O}}\right\rangle _{V,l+1/2}^{0}  &  =\frac
{2\pi}{m_{e}}\frac{\widetilde{q}_{l+1/2}}{\widehat{q}_{l+1/2}}\sum
_{i=1}^{n_{p}}\sum_{j=0}^{n_{\xi_{0}}-1}\frac{p_{i}^{3}}{\gamma_{i}}%
\lambda^{l+1/2,j+1/2}S_{p,l+1/2,i,j+1/2}^{\left(  0\right)  \mathcal{O}%
}\left(  f_{0}^{\left(  0\right)  }\right)  \Delta p_{i}\Delta\xi
_{0,j+1/2}\label{pabsdis}\\
\left\langle P_{abs}^{\mathcal{O}}\right\rangle _{V,l+1/2}^{1}  &  =\frac
{2\pi}{m_{e}}\frac{\widetilde{q}_{l+1/2}}{\widehat{q}_{l+1/2}}\sum
_{i=1}^{n_{p}}\sum_{j=0}^{n_{\xi_{0}}-1}\frac{p_{i}^{3}}{\gamma_{i}}%
\lambda^{l+1/2,j+1/2}S_{p,l+1/2,i,j+1/2}^{\left(  0\right)  \mathcal{O}%
}\left(  g^{\left(  0\right)  }\right)  \Delta p_{i}\Delta\xi_{0,j+1/2}\\
\left\langle \widetilde{P}_{abs}^{\mathcal{O}}\right\rangle _{V,l+1/2}^{1}  &
=\frac{2\pi}{m_{e}}\frac{\widetilde{q}_{l+1/2}}{\widehat{q}_{l+1/2}}\sum
_{i=1}^{n_{p}}\sum_{j=0}^{n_{\xi_{0}}-1}\frac{p_{i}^{3}}{\gamma_{i}}%
\lambda^{l+1/2,j+1/2}\widetilde{S}_{p,l+1/2,i,j+1/2}^{\left(  0\right)
\mathcal{O}}\left(  \widetilde{f}^{\left(  0\right)  }\right)  \Delta
p_{i}\Delta\xi_{0,j+1/2}%
\end{align}

The discretization of the momentum-space flux components $S_{p,l+1/2,i,j+1/2}%
^{\left(  0\right)  \mathcal{O}}$ and $\widetilde{S}_{p,l+1/2,i,j+1/2}%
^{\left(  0\right)  \mathcal{O}}$ is done in (\ref{sp0dis}).

\subsubsection{Stream Function for Momentum Space fluxes}

\bigskip The stream function gives the local direction of the momentum-space
fluxes, and its gradient is an indication of the flux intensity. Because it is
a flux function, it is naturally defined on the momentum-space flux grid.
According to the three equivalent expressions for the stream function
(\ref{stream1}-\ref{stream2}), we have the following discretizations, for
$1\leq i\leq n_{p}$ and $1\leq j\leq n_{\xi}-1$
\begin{align}
A_{l+1/2,i,j}^{\left(  0\right)  }  &  =\frac{2\pi p_{i}^{2}}{n_{e,l+1/2}}%
\sum_{m=0}^{j-1}\Delta\xi_{0,m+1/2}\;S_{p,l+1/2,i,m+1/2}^{\left(  0\right)
}\label{streamdis}\\
A_{l+1/2,i,j}^{\left(  0\right)  }  &  =\frac{2\pi p_{i}^{2}}{n_{e,l+1/2}}%
\sum_{m=j}^{n_{\xi}-1}\Delta\xi_{0,m+1/2}\;S_{p,l+1/2,i,m+1/2}^{\left(
0\right)  }%
\end{align}
and
\begin{equation}
A_{l+1/2,i,j}^{\left(  0\right)  }=\frac{2\pi\sqrt{1-\xi_{0,j}^{2}}%
}{n_{e,l+1/2}}\sum_{m=0}^{i-1}p_{m+1/2}\Delta p_{m+1/2}\;S_{\xi,l+1/2,m+1/2,j}%
^{\left(  0\right)  } \label{streamdis2}%
\end{equation}
where the boundary conditions are%
\begin{equation}
A_{l+1/2,0,j}^{\left(  0\right)  }=A_{l+1/2,i,0}^{\left(  0\right)
}=A_{l+1/2,i,n_{\xi}}^{\left(  0\right)  }=0
\end{equation}

The discretization of the momentum-space flux components $S_{p,l+1/2,i,j+1/2}%
^{\left(  0\right)  \mathcal{O}}$ and $S_{\xi,l+1/2,i+1/2,j}^{\left(
0\right)  \mathcal{O}}$ is done in (\ref{sp0dis}) and (\ref{sxi0dis}).

\subsubsection{Ohmic electric field}

As shown in Sec. \ref{Efield}, the flux surface averaged parallel Ohmic
electrid field $\left\langle E_{\Vert}\right\rangle _{\phi}\left(
\psi\right)  $ may be expressed as a function of its local value $E_{\Vert
0}\left(  \psi\right)  $ taken at the poloidal position where the magnetic
field $\mathbf{B}$ is minimum, and
\begin{equation}
\left\langle E_{\Vert}\right\rangle _{\phi,l+1/2}=E_{\Vert0,l+1/2}%
\frac{\widetilde{q}_{l+1/2}}{\overline{q}_{l+1/2}}\frac{R_{p}}{R_{0,l+1/2}%
}\frac{B_{T0,l+1/2}}{B_{0,l+1/2}}\overline{\lambda}_{1,-3,4}^{l+1/2,j+1/2}%
\end{equation}

\subsubsection{Exact and effective fractions of trapped electrons}

The exact fraction of trapped electrons is given by relation,
\begin{align}
&\mathcal{F}_{t,l+1/2}=\left[\sum_{i=0}^{n_{p}-1}\sum_{j=0}^{n_{\xi_{0}}%
-1}p_{i+1/2}^{2}H\left(  \left|  \xi_{0,j+1/2}\right|  -\xi_{0T,l+1/2}\right)
\lambda^{l+1/2,j+1/2}\right.\nonumber\\&\times\left.\left(  \widetilde{f}_{l+1/2,i+1/2,j+1/2}^{(0)}%
+g_{l+1/2,i+1/2,j+1/2}^{(0)}\right)  \Delta
p_{i+1/2}\Delta\xi_{0,j+1/2}\right]\nonumber\\&\times\left[
\sum_{i=0}^{n_{p}-1}\sum_{j=0}^{n_{\xi_{0}}-1}p_{i+1/2}^{2}\lambda
^{l+1/2,j+1/2}\right.\nonumber\\&\times\left.\left(  f_{0,l+1/2,i+1/2,j+1/2}^{(0)}+\widetilde{f}%
_{l+1/2,i+1/2,j+1/2}^{(0)}+g_{l+1/2,i+1/2,j+1/2}^{(0)}\right)
\Delta p_{i+1/2}\Delta\xi_{0,j+1/2}\right]^{-1}\nonumber\\%
\end{align}
according to calculations given in Sec. \ref{moments}.

The effective fraction of trapped electrons, as deduced from the Lorentz model
in Sec. \ref{dkediscretinit} is%

\begin{align}
\mathcal{F}_{t,l+1/2}^{eff.}  &  =\frac{3}{2}\frac{B_{T0,l+1/2}}{B_{0,l+1/2}%
}\times\nonumber\\
&  \left[  \frac{\widetilde{q}_{l+1/2}}{\overline{q}_{l+1/2}}\frac
{B_{T0,l+1/2}}{B_{0,l+1/2}}\sum_{j=0}^{n_{\xi_{0}}-1}\xi_{0,j+1/2}^{2}%
\lambda_{2,-2,2}^{l+1/2,j+1/2}\Delta\xi_{0,j+1/2}\right. \nonumber\\
&  \left.  -\frac{q_{l+1/2}}{\overline{q}_{l+1/2}}\frac{R_{0,l+1/2}}{R_{p}%
}\sum_{j=0}^{n_{\xi_{0}}-1}\sigma_{j+1/2}H\left(  \left|  \xi_{0,j+1/2}%
\right|  -\xi_{0T,l+1/2}\right) \right.\nonumber\\&\times\left.
\;\xi_{0,j+1/2}I_{\mathcal{L}}\left(
\psi_{l+1/2},\left|  \xi_{0,j+1/2}\right|  \right)  \Delta\xi_{0,j+1/2}%
\right]
\end{align}

\subsubsection{Runaway loss rate}

As shown in Sec. \ref{moments}, the runaway loss rate $\left\langle \Gamma
_{R}\right\rangle _{V}\left(  \psi\right)  $ is given%

\begin{equation}
\left\langle \Gamma_{R}\right\rangle _{V,l+1/2}=\frac{\widetilde{q}_{l+1/2}%
}{\widehat{q}_{l+1/2}}2\pi p_{n_{p}-1}^{2}\sum_{j=0}^{n_{\xi_{0}}-1}%
\lambda^{l+1/2,j+1/2}S_{p,l+1/2,n_{p}-1,j+1/2}^{\left(  0\right)  }\Delta
\xi_{0,j+1/2}%
\end{equation}

\subsubsection{Magnetic ripple losses}

In a similar way, the magnetic ripple loss rate $\Gamma_{ST}^{\left(
0\right)  }\left(  \psi\right)  $ is given by
\begin{align}
\Gamma_{ST,l+1/2}^{\left(  0\right)  }  &  =2\pi p_{i_{c}}^{2}\sum
_{j=0}^{n_{\xi_{0}}-1}\lambda^{l+1/2,j+1/2}\left(  1-H\left(  \left|
\xi_{0,j+1/2}\right|  -\xi_{0ST,l+1/2}\right)  \right)  S_{p,l+1/2,i_{c}%
,j+1/2}^{\left(  0\right)  }\Delta\xi_{0,j+1/2}\nonumber\\
&  +4\pi\lambda^{l+1/2,j_{ST}}\sqrt{1-\xi_{0ST,j}^{2}}\sum_{i=0}^{n_{p}%
-1}p_{i+1/2}H\left(  p_{i+1/2}-p_{i_{c}+1/2}\right)  S_{\xi_{0}%
,l+1/2,i+1/2,j_{ST}}^{\left(  0\right)  }\Delta p_{i+1/2}%
\end{align}
and the second formulation given in Sec.\ref{moments} is
\begin{align}
\Gamma_{ST,l+1/2}^{\left(  0\right)  }  &  =2\pi\sum_{i=0}^{n_{p}-1}\sum
_{j=0}^{n_{\xi_{0}}-1}\left(  1-H\left(  \left|  \xi_{0,j+1/2}\right|
-\xi_{0ST,l+1/2}\right)  \right)  H\left(  p_{i+1/2}-p_{i_{c}+1/2}\right)
\nonumber\\
&  \times\nu_{d_{ST,l+1/2,i+1/2,j+1/2}}f_{0,l+1/2,i+1/2,j+1/2}^{(0)}\Delta
\xi_{0,j+1/2}\Delta p_{i+1/2}%
\end{align}
where $i_{c}$ is the index number that corresponds to the detrapping threshold
by collisions, and $\xi_{0ST,l+1/2}$ is the picth-angle boundary value between
super-trapped and trapped electrons.

\subsubsection{Non-thermal bremsstrahlung}

For the bremsstrahlung emission, it is necessary to calculate numericaly all
coefficients of the Legendre series. Since Legendre polynomials $P_{m}$
strongly oscillate between $-1$ and $+1$, as their order $m$ increases, it is
not possible to evaluate accurately integrals of the type
\begin{equation}
h^{\left(  m\right)  }=\int_{-1}^{+1}h\left(  x\right)  P_{m}\left(  x\right)
dx \label{projLegendre}%
\end{equation}
by standard integration techniques, like trapezoidal or Simpson rules. The
only possibility is to replace integral (\ref{projLegendre}) by a discrete
sum, namely a Gaussian quadrature,
\begin{equation}
h^{\left(  m\right)  }=\sum_{n=1}^{N}w_{n}h\left(  x_{n}\right)  P_{m}\left(
x_{n}\right)
\end{equation}
where weights $w_{n}$ and abscissas $x_{n}$ are determined independently.
Here, the use of Legendre polynomials lead to consider the fast ans accurate
Gauss-Legendre algorithm as derived by G. B. Ribicki in Ref. \cite{pre92},
where $x_{n}$ are $N$ zeros of the Legendre polynomial of degree $N$, the
weights being given by the relation
\begin{equation}
w_{n}=\frac{2}{\left(  1-x_{n}^{2}\right)  \left[  P_{N}^{\prime}\left(
x_{n}\right)  \right]  ^{2}}%
\end{equation}

Very accurate determination of $h^{\left(  m\right)  }$ may be obtained by
this method, which requires a value of $N=\;50$.
